% Paper 3A manuscript. Built and checked by paper/tex/build_ms.py (gates G1-G7). Every number in the text is
% \V{key}{value} bound to a results file or a frozen document in paper/tex/bindings.json.
% Sources: paper/manuscript/introduction_draft.md and methods_draft.md (2026-11-09), PREREG_P3A.md (frozen at
% 53cbc39, tag prereg-p3a-freeze). The journal template is applied after the venue is chosen (hub task T7).
\documentclass[11pt]{article}
\usepackage[a4paper,margin=25mm]{geometry}
\usepackage[T1]{fontenc}
\usepackage{lmodern}
\usepackage{amsmath,amssymb}
\usepackage{graphicx}
\usepackage{framed}
\usepackage{booktabs}
\usepackage{textcomp}
\usepackage[hidelinks]{hyperref}
% Macros shared by main.tex and si.tex. The build gates (paper/tex/build_ms.py) read these names.
% \V{key}{value}  a number (math mode); key -> paper/tex/bindings.json; G1 compares value with the source
% \W{key}{word}   a count written as a word (text); same binding rules (fmt word / Word)
% \C{...}         a mathematical constant inside a formula (2 in 2 pi, 0 in p -> 0); no data
% \ID{...}        an identifier that contains digits (hash, date, version, DOI, ORCID, model-cell names); no data
% \snote{n}       a Supplementary Note number; G2 checks that si.tex defines it
% \pending{name}{text}  a result not yet available; G7 lists it and a --final build refuses it
\newcommand{\V}[2]{\ensuremath{#2}}
\newcommand{\W}[2]{#2}
\newcommand{\C}[1]{#1}
\newcommand{\ID}[1]{#1}
\newcommand{\snote}[1]{#1}
\newcommand{\pending}[2]{\par\noindent\fbox{\parbox{0.95\linewidth}{\textit{Pending result (#1).} #2}}\par}
\newcommand{\pss}{p/p_{\mathrm{sat}}}
\newcommand{\qref}{q_{\mathrm{ref}}}
\newcommand{\Qst}{Q_{\mathrm{st}}}


\title{Hard thermodynamic constraints and the temperature extrapolation of learned water isotherms in
metal-organic and covalent organic frameworks: a pre-registered test}
\author{Abdulhamid Batayhi\\ Yildiz Technical University, Istanbul, Turkey}
\date{}

\begin{document}
\maketitle

\begin{abstract}
\pending{abstract}{Written from the confirmatory verdicts (results/verdicts\_confirmatory.json) after the analysis.}
\end{abstract}

\section*{Introduction}

Water sorbents built from metal-organic frameworks (MOFs) and covalent organic frameworks (COFs) are studied for
atmospheric water harvesting\cite{solovyeva2021,nguyen2020} and for adsorption heat pumps and
chillers\cite{henninger2012,lenzen2019}. These devices run between temperatures: a MOF-801 cooling cycle adsorbs at
\V{c34_ads}{30}~\textdegree C and regenerates at \V{c34_reg_lo}{80} to \V{c34_reg_hi}{85}~\textdegree
C\cite{solovyeva2018}. Its design needs the uptake at temperatures where no isotherm was measured, and the isosteric
heat, which enters the cycle model\cite{lenzen2019}. Both rest on few temperatures. A practical guide asks for
isotherms spaced by \V{c06_dT}{20}~K or less for a Clausius-Clapeyron heat and advises against extrapolation beyond
the measured range\cite{nuhnen2020}. For Mg-SU-102, isotherms were measured from \V{c33_lo}{15} to
\V{c33_hi}{55}~\textdegree C, but only those from \V{c33_heat_lo}{25} to \V{c33_heat_hi}{40}~\textdegree C had
enough points for the heat\cite{oppenheim2025}.

The standard temperature model for water in MOFs is the characteristic curve: the uptake is written as a function
of the adsorption potential $A = RT\ln(p_{\mathrm{sat}}/p)$, the isotherms of all temperatures are expected to fall
on one curve, and that curve is used to model cycles\cite{liu2020}. Within this approach the isosteric heat is the
vaporisation enthalpy plus $A$\cite{lenzen2019}. Invariance is usually shown by overlaying curves, without an error
number, as for MIL-160 from \V{c14_lo}{20} to \V{c14_hi}{50}~\textdegree C\cite{solovyeva2021} and Al-fumarate
from \V{c35_lo}{298} to \V{c35_hi}{333}~K\cite{trezza2022}. The Mg-SU-102 curves shift slightly towards lower
energies as the temperature rises\cite{oppenheim2025}, and MOF-303 loses a large part of its capacity between
\V{c59_lo}{25} and \V{c59_hi}{45}~\textdegree C\cite{hastings2024}. Statistical thermodynamics gives a
characteristic curve when added sorbate contributes purely energetically, and a temperature-dependent isotherm
when sorbate number correlates with entropy\cite{shimizu2021}. The other common route, Langmuir or Sips sites with
temperature-dependent affinities fitted to all temperatures at once\cite{lenzen2019,siperstein2021}, has its own
defects. Sips and Freundlich terms have no finite slope at zero pressure and their heats diverge at low coverage,
dual-site Langmuir heats can show an unphysical maximum, and models that fit equally well disagree on the heat
wherever they extrapolate\cite{siperstein2021}.

Machine-learned isotherm models add flexibility but little thermodynamics. In two recent examples temperature is a
plain input: one network predicted methane isotherms of \V{l24_n}{17{,}644} MOFs from \V{l24_lo}{278} to
\V{l24_hi}{338}~K and was extrapolated to \V{l24_x1}{268} and \V{l24_x2}{358}~K\cite{dong2025}; another predicts
water isotherms from pressure and structural properties\cite{liu2023}. Where heats appear, they are separate
outputs. A neural ordinary differential equation learns $\mathrm{d}q/\mathrm{d}\ln p$ at one temperature, observes
monotonicity without imposing it, and predicts the heat with a second network\cite{lin2025}. A zeolite study joined
a learned loading at \V{l04_T}{300}~K and an independently learned heat through an integrated Clausius-Clapeyron
step with a temperature-independent heat\cite{daou2024}. When physics enters, it enters as a penalty. A network
for hydrogen sorption adds monotonicity as a weighted loss term and reports \V{l02_sat}{98.6}~\% monotonicity
satisfaction, so \V{l02_viol}{1.4}~\% of held-out cases violate it\cite{nooraiepour2026}. A penalty-trained network
is not provably monotone; monotonicity can only be certified afterwards\cite{liu2020cert}.

A different construction is possible. The grand (immersion) potential of the adsorbed phase,
$\Omega = -RT\int n\,\mathrm{d}\ln p$, has the isotherm as its derivative with respect to $\ln
p$\cite{myers2002,coudert2008}. A model can learn this one scalar as a function of $\ln p$ and $T$. If it is convex
in $\ln p$, the isotherm is monotone for every parameter value, and input-convex networks give convexity by
construction\cite{amos2017}. A potential linear in $p$ at low pressure gives a finite Henry slope, which Sips terms
lack\cite{siperstein2021}. The isosteric heat then follows by Clausius-Clapeyron from second derivatives of the
same function and is defined wherever $\mathrm{d}q/\mathrm{d}\ln p$ is positive. Neighbouring fields already learn
potentials this way. A learned free energy gives fluid properties as its derivatives, with exact Maxwell
relations\cite{rosenberger2022}. In mixture thermodynamics, a graph network for the excess Gibbs energy replaced a
Gibbs-Duhem penalty\cite{rittig2023} and needs no extra loss terms\cite{rittig2024}, and hard-constraint networks
build the Gibbs-Duhem relation into the architecture\cite{specht2024} and take the excess enthalpy from the
temperature derivative of a learned excess Gibbs energy\cite{hoffmann2026}. In adsorption, a learned density
functional has produced an isotherm of a Lennard-Jones fluid in a model pore\cite{cheng2026}. In the sources we
searched, no model learns such a potential from measured isotherms of real sorbents or tests it on water.
% lit_3a_ml.md section F, with its caveats: "not found in the sources searched"; ChemRxiv and the Liu 2023 full text are blind spots

Hard constraints make a model admissible, but not necessarily accurate. For phase-transition enthalpies, a hard
thermodynamic constraint raised the error by more than \V{l47_err}{10}~kJ\,mol$^{-1}$, while a soft one kept the
accuracy\cite{leenhouts2025}. Two parts of the construction can also act separately: the shape constraint fixes
monotonicity and the Henry limit, while the temperature structure (a van't Hoff affinity per term, a free
dependence on $\C{1}/T$, or a characteristic curve) decides how the isotherm moves between temperatures. The open
question is which of the two, if either, makes extrapolation in temperature better than penalties, generic networks
and classical models.

The present study therefore asks whether hard constraints on a learned adsorbed-phase potential make its
temperature extrapolation of measured water isotherms better than penalties, generic networks and classical models.
It does so in a pre-registered factorial design that crosses the shape constraint (none, penalty or hard) with the
temperature structure (free in $\C{1}/T$, van't Hoff per term, or a temperature-invariant characteristic curve). The
learned potential is the cell with hard constraints and van't Hoff terms. It is compared with its penalty and
unconstrained twins, the hard-constrained cells of the other two temperature structures, a generic neural network
and a dual-site Sips model with van't Hoff affinities. The data are \V{n_conf}{13} real sets of water isotherms
measured at \W{min_T_word}{three} or more temperatures, \V{n_mof}{5} on MOFs and \V{n_cof}{8} on COFs. Each model
predicts a material's highest and lowest temperatures from the others. The design, tests and words of every
verdict were fixed before any confirmatory fit, including words for a classical model that wins and for models
equivalent within \V{sesoi_pct}{25}~\%. Every defect found in the design, code or data is recorded in a public
correction ledger. None of the \V{n_conf}{13} materials has a calorimetric water heat; our search found one
verified value for a MOF with multi-temperature isotherms, a single integral heat for MOF-801\cite{kim2016}, and
none for any COF. We therefore test the heat on synthetic truths with exact heats, built on each material's own
temperatures and pressures.
\pending{contributions}{Numbered contributions, each with its headline number, evidence label and section pointer
(lessons file, addendum on the pinn-machining style guide), written from the confirmatory verdicts.}

\section*{Results}

\pending{results-design}{Figure~1 (paper/figures/fig\_design.pdf): the factorial, the tasks and the evidence label
of every comparison. Fits completed and failed per cell, from results/confirmatory\_runs.shard*.jsonl.}

\pending{results-H1}{H1b (primary, non-inferiority) to H1e: R with exact 95~\% and 90~\% intervals, Holm-adjusted
p and verdict words, from results/verdicts\_confirmatory.json (paper/table\_verdicts.py); Figure~2a forest
(paper/fig\_verdicts.py). H1a descriptive.}

\pending{results-H3}{H3 admissibility: violating fits per cell over the E-high and E-low fits, Wilson intervals,
the pre-registered prediction (V-pen any class but none and P none); Figure~2b. The P bug check (P\_bug\_fits).}

\pending{results-H2}{H2 isosteric heat on the synthetic truths (results/synthetic\_h2 confirmatory runs), the 15
tests of PREREG 7.1.}

\pending{results-descriptive}{Descriptive analyses fixed in PREREG 6.4: I-mid, seed spread, the MLP, F-rel, P-K,
subgroups; the sensitivity variants (scripts/sensitivity.py verdicts NAME, seven variants) and MOF-801.}

\section*{Discussion}

\pending{discussion}{What the verdicts mean for hard constraints against penalties (accuracy against
admissibility), the temperature structure, limitations (n = 13, digitised data, no calorimetric heats, two
training temperatures in E tasks for most materials) and what is not claimed (PREREG section 9).}

\section*{Methods}

\subsection*{Pre-registration}

The design, the data, the models, the tuning rules, the metrics, the statistical tests and the words of every
verdict were fixed in a pre-registration and frozen before any confirmatory fit. The freeze is a public, tagged
commit (\texttt{prereg-p3a-freeze}) whose manifest holds the SHA-256 of \V{n_frozen}{71} files (the text, the data,
the tuned configurations and the fits behind them, the synthetic-world and power results, and every script that
builds, fits or scores, including an independently written second scorer) and the package versions; the
confirmatory scripts refuse to run if any of them differs. Every defect found in the design, code or data, before or
after the freeze, is listed with its cause and fix in a public ledger.
% CORRECTIONS_P3A.md; P3A-1 to P3A-25 at the freeze; update at submission

\subsection*{Data}

Water adsorption isotherms of MOFs and COFs measured at several temperatures on one sample were collected from the
primary literature and from the NIST adsorption database. A material entered the confirmatory set if, at
\W{rule1_word}{three} or more temperatures, its adsorption isotherm had at least \V{rule2_n}{8} points with
$\C{0} < \pss \le \V{rule2_cut}{0.90}$, its figure passed the digitisation quality gates, its temperatures were
unambiguous, and its curve was covered: no gap above \V{rule5_gap}{0.25} in $\pss$, and a first point at
$\pss \le \V{rule5_start}{0.10}$ or at an uptake of at most \V{rule5_q}{0.10}~$\qref$. Where a paper gave both
measured isotherms and curves derived from them, the measured isotherms were used. The confirmatory set has
\V{n_conf}{13} materials, \V{n_mof}{5} MOFs and \V{n_cof}{8} COFs; every decision and its reason is in the inclusion
log. \W{n_dev_word}{Four} further materials form a development set used only to choose hyperparameters.

Isotherms published only as figures were digitised before any model was fitted. Vector figures were read from the
PDF drawing commands, with marker centres mapped through the axis ticks; raster figures by template matching with
the figure's own marker footprint, with tick marks located to sub-pixel precision and an open-marker veto against
desorption points. No physical prior was used to filter points. The method was validated against published numeric
data. On Mg-SU-102 the vector extraction recovered \V{dig_vec_n}{198} of \V{dig_vec_truth}{198} points with a
root-mean-square uptake error of $\V{dig_vec_rms}{7\times10^{-6}}$~g\,g$^{-1}$. On raster figures of HKUST-1 the
root-mean-square uptake error was \V{dig_ras_lo}{0.11} to \V{dig_ras_hi}{0.19}~\% of the axis span, or
\V{dig_rel_lo}{0.0012} to \V{dig_rel_hi}{0.0034} of the isotherm's maximum, with no false detections.

Only adsorption branches with $\pss \le \V{rule2_cut}{0.90}$ were used. Relative pressures were converted with the
IAPWS saturation pressure, absolute pressures were used as reported, and uptakes were converted to
mol\,kg$^{-1}$. The scale $\qref$ of a material, its largest uptake in this range over all temperatures, is used
only for scoring.

\subsection*{Models}

Every model is a function $q(s,T)$ of $s = \ln(p/\mathrm{Pa})$ and temperature, and the isosteric heat of every
model is computed by the same code, $\Qst = -RT^{2}(\partial q/\partial T)/(\partial q/\partial s)$, by automatic
differentiation.

The potential model (P) learns the reduced spreading pressure
\begin{equation}
\psi(s,T) = \sum_k \frac{c_k(T)}{g_k}\,\mathrm{softplus}\bigl(g_k\,(s-\theta_k(T))\bigr),
\qquad q^{*} = \frac{\partial\psi}{\partial s} = \sum_k c_k(T)\,\sigma\bigl(g_k\,(s-\theta_k(T))\bigr),
\end{equation}
with capacities $c_k > \C{0}$, Henry terms ($g_k = \C{1}$) and cooperative terms ($g_k \ge \V{gcoop}{1.5}$), van't
Hoff positions $\theta_k(T) = \theta_{k0} - (Q_k/R)(\C{1}/T - \C{1}/T_{\mathrm{ref}})$ and an adsorbate-expansion
factor. For every parameter value, $q^{*}$ increases strictly with $p$, $q^{*}/p$ tends to a positive Henry constant as
$p \to \C{0}$, and
the heat is defined everywhere.

The models form a factorial of shape constraint (none, penalty, hard) and temperature structure: a free network in
$\C{1}/T$ (F), van't Hoff per term (V), or the temperature-invariant characteristic curve in $A$ (A). In the V and A
columns, the unconstrained and penalised models have the function class of P or of the hard characteristic curve
without the constraints, the same initial function and the same tuned structure; the penalty models add
monotonicity and (except for the characteristic curve) a scale-free Henry-limit penalty on collocation points that
reach the held-out temperature. A dual-site Sips model with van't Hoff affinities (C1) is the classical comparator.
A multilayer perceptron on $(\ln \pss, \C{1}/T)$ and a Kirchhoff variant of P with a heat capacity per term are reported
descriptively.

\subsection*{Fitting and tuning}

Each model was fitted by Adam followed by L-BFGS on the mean squared error of $q$ scaled by the largest training
uptake, in double precision, from \V{n_seeds}{5} seeds; the fit with the lowest training loss was kept, never
selected on held-out error. Every scale of a fit came from the training temperatures only. Hyperparameters were
chosen on the development set alone, in two stages: the models that own a structure tried up to
\V{n_cfg}{12} configurations each, and the unconstrained and penalised models inherited the chosen structure, the
penalty models tuning only the penalty weight over $\V{lam_lo}{10^{-4}}$ to $\V{lam_hi}{100}$.

\subsection*{Tasks and outcomes}

For a material with temperatures $T_1 < \dots < T_k$, the task E-high trains on $T_1,\dots,T_{k-1}$ and predicts
$T_k$; E-low trains on $T_2,\dots,T_k$ and predicts $T_1$. The primary outcome is the held-out error
$E = \mathrm{RMSE}(q_{\mathrm{pred}} - q_{\mathrm{obs}})/\qref$, summarised per material as the mean of
$\log_{10} E$ over the two tasks. Physical violations were counted on a dense grid: a fit violated if an isotherm
fell by more than \V{viol_pct}{2}~\% of $\qref$ below its running maximum, or went below
$-\V{viol_pct}{2}$~\% of $\qref$, inside the domain where the penalties act; violations were counted over the fits
of the two extrapolation tasks. Because no material here has a calorimetric water heat, the heat was tested on
synthetic isotherms with exact heats, built on every material's own temperatures and pressures from
\W{n_truth_word}{five} truth families, with noise of \V{sigma_pct}{0.4}~\% of $\qref$.

\subsection*{Statistics}

The unit is the material ($n = \V{n_conf}{13}$). For each comparison of P with a model X,
$d_m = e_m(\mathrm{P}) - e_m(\mathrm{X})$ was tested with the exact two-sided Wilcoxon signed-rank test, with Holm
correction over the \W{n_holm_word}{four} confirmatory comparisons (penalty model, characteristic curve, generic
network, dual-site Sips). The effect is the Hodges-Lehmann estimate with exact distribution-free
\V{ci_main}{95}~\% and \V{ci_aux}{90}~\% intervals, reported as the error ratio $R = \C{10}^{d}$. A failed fit was
ranked worst, never dropped. The words of each verdict follow a table fixed in advance, with a smallest effect of
interest of \V{sesoi_pct}{25}~\% ($R = \V{sesoi_lo}{0.8}$ or $\V{sesoi_hi}{1.25}$). The one declared primary claim,
that the hard constraints cost no more than \V{sesoi_pct}{25}~\% in accuracy against the penalty model, was judged
by non-inferiority: the upper exact \V{ci_main}{95}~\% bound of $R$ below \V{sesoi_hi}{1.25} (one-sided
\V{ni_alpha}{2.5}~\%), with the \V{ci_holm}{98.75}~\% bound reported beside it. At \V{n_conf}{13} materials its exact
size at the margin is \V{size95}{0.024} (\V{size9875}{0.005} for the \V{ci_holm}{98.75}~\% bound). Before the
freeze, the power of each verdict was simulated with the analysis code itself, and the whole pipeline was run on
\W{n_worlds_word}{five} synthetic worlds, with an expected outcome written for each before it ran, to show which
verdicts can fail (Supplementary Note~\snote{1}). \W{n_missed_word}{Three} of the expectations were missed and are
reported as missed, and the comparison of a free with a van't Hoff temperature dependence was made descriptive
because it could not fail.

\section*{Data Availability}

The pre-registration, the freeze manifest, the corrections ledger, the digitisation provenance and all results are
public at \url{https://github.com/abdulhamidbatayhi123/ai-isotherm-thermo}; an archived release with a DOI will be
cited. Digitised point lists from sources whose licence does not permit redistribution are kept locally; their
provenance, fitted parameters and predictions are published.

\section*{Code Availability}

All code is in the same repository and archive.

\section*{Use of AI tools}

\pending{ai-statement}{The author's saved wording (AI as a tool; the author designed the study, made every
decision and takes responsibility), adapted to the target journal's policy.}

% written by build_ms.py from refs_master.json (order of first citation); do not edit
\renewcommand\refname{References}
\begin{thebibliography}{99}
\bibitem{solovyeva2021} Solovyeva, M., Krivosheeva, I., Gordeeva, L. \& Aristov, Y. MIL-160 as an adsorbent for atmospheric water harvesting. \textit{Energies} \textbf{14}, 3586 (2021). doi:10.3390/en14123586
\bibitem{nguyen2020} Nguyen, H. L., Hanikel, N., Lyle, S. J., Zhu, C., Proserpio, D. M. \& Yaghi, O. M. A porous covalent organic framework with voided square grid topology for atmospheric water harvesting. \textit{J. Am. Chem. Soc.} \textbf{142}, 2218 (2020). doi:10.1021/jacs.9b13094
\bibitem{henninger2012} Henninger, S. K., Jeremias, F., Kummer, H. \& Janiak, C. MOFs for use in adsorption heat pump processes. \textit{Eur. J. Inorg. Chem.} \textbf{2012}, 2625 (2012). doi:10.1002/ejic.201101056
\bibitem{lenzen2019} Lenzen, D. et al. A metal-organic framework for efficient water-based ultra-low-temperature-driven cooling. \textit{Nat. Commun.} \textbf{10}, 3025 (2019). doi:10.1038/s41467-019-10960-0
\bibitem{solovyeva2018} Solovyeva, M. V., Gordeeva, L. G., Krieger, T. A. \& Aristov, Yu. I. MOF-801 as a promising material for adsorption cooling: equilibrium and dynamics of water adsorption. \textit{Energy Convers. Manag.} \textbf{174}, 356 (2018). doi:10.1016/j.enconman.2018.08.032
\bibitem{nuhnen2020} Nuhnen, A. \& Janiak, C. A practical guide to calculate the isosteric heat/enthalpy of adsorption via adsorption isotherms in metal-organic frameworks, MOFs. \textit{Dalton Trans.} \textbf{49}, 10295 (2020). doi:10.1039/d0dt01784a
\bibitem{oppenheim2025} Oppenheim, J., Yang, Z., Dinakar, B. \& Dinc\u{a}, M. High-capacity water sorbent cycles without hysteresis under dry conditions. \textit{Nat. Commun.} \textbf{16}, 4297 (2025). doi:10.1038/s41467-025-59551-2
\bibitem{liu2020} Liu, X., Wang, X. \& Kapteijn, F. Water and metal-organic frameworks: from interaction toward utilization. \textit{Chem. Rev.} \textbf{120}, 8303 (2020). doi:10.1021/acs.chemrev.9b00746
\bibitem{trezza2022} Trezza, G., Bergamasco, L., Fasano, M. \& Chiavazzo, E. Minimal crystallographic descriptors of sorption properties in hypothetical MOFs and role in sequential learning optimization. \textit{npj Comput. Mater.} \textbf{8}, 123 (2022). doi:10.1038/s41524-022-00806-7
\bibitem{hastings2024} Hastings, J. et al. High-temperature water adsorption isotherms and ambient temperature water diffusion rates on water harvesting metal-organic frameworks. \textit{J. Phys. Chem. C} \textbf{128}, 11328 (2024). doi:10.1021/acs.jpcc.4c01733
\bibitem{shimizu2021} Shimizu, S. \& Matubayasi, N. Temperature dependence of sorption. \textit{Langmuir} \textbf{37}, 11008 (2021). doi:10.1021/acs.langmuir.1c01576
\bibitem{siperstein2021} Siperstein, F. R., Avenda\~{n}o, C., Ortiz, J. J. \& Gil-Villegas, A. Analytic expressions for the isosteric heat of adsorption from adsorption isotherm models and two-dimensional SAFT-VR equation of state. \textit{AIChE J.} \textbf{67}, e17186 (2021). doi:10.1002/aic.17186
\bibitem{dong2025} Dong, M.-E., Wu, X. \& Cai, W. Prediction of methane adsorption isotherms in metal-organic frameworks by neural networks: two-dimensional energy gradient feature and masked learning mechanism. \textit{J. Phys. Chem. B} \textbf{129}, 11333 (2025). doi:10.1021/acs.jpcb.5c05384
\bibitem{liu2023} Liu, Z., Shen, D., Cai, S., Tu, Z. \& Li, S. Machine learning-assisted prediction of water adsorption isotherms and cooling performance. \textit{J. Mater. Chem. A} \textbf{11}, 19455 (2023). doi:10.1039/d3ta03586g
\bibitem{lin2025} Lin, E., Wang, E. \& Deng, S. Neural ordinary differential equations (ODEs) for smooth, high-accuracy isotherm reconstruction, interpolation, and extrapolation. \textit{npj Comput. Mater.} \textbf{11}, 348 (2025). doi:10.1038/s41524-025-01819-8
\bibitem{daou2024} Daou, A. S. S., Fang, H., Boulfelfel, S. E., Ravikovitch, P. I. \& Sholl, D. S. Machine learning and IAST-aided high-throughput screening of cationic and silica zeolites for alkane capture, storage, and separations. \textit{J. Phys. Chem. C} \textbf{128}, 6089 (2024). doi:10.1021/acs.jpcc.4c00066
\bibitem{nooraiepour2026} Nooraiepour, M., Masoudi, M., Song, Z. \& Hellevang, H. Thermodynamically constrained deep learning for hydrogen sorption prediction: overcoming classical isotherm generalization failure through physics-informed cross-lithology transferability. \textit{ACS Omega} (2026). doi:10.1021/acsomega.6c08678
\bibitem{liu2020cert} Liu, X., Han, X., Zhang, N. \& Liu, Q. Certified monotonic neural networks. In \textit{Advances in Neural Information Processing Systems 33}, 15427 (2020). Preprint at https://arxiv.org/abs/2011.10219
\bibitem{myers2002} Myers, A. L. Thermodynamics of adsorption in porous materials. \textit{AIChE J.} \textbf{48}, 145 (2002). doi:10.1002/aic.690480115
\bibitem{coudert2008} Coudert, F.-X., Jeffroy, M., Fuchs, A. H., Boutin, A. \& Mellot-Draznieks, C. Thermodynamics of guest-induced structural transitions in hybrid organic-inorganic frameworks. \textit{J. Am. Chem. Soc.} \textbf{130}, 14294 (2008). doi:10.1021/ja805129c
\bibitem{amos2017} Amos, B., Xu, L. \& Kolter, J. Z. Input convex neural networks. In \textit{Proc. 34th International Conference on Machine Learning}, PMLR 70, 146 (2017). Preprint at https://arxiv.org/abs/1609.07152
\bibitem{rosenberger2022} Rosenberger, D., Barros, K., Germann, T. C. \& Lubbers, N. Machine learning of consistent thermodynamic models using automatic differentiation. \textit{Phys. Rev. E} \textbf{105}, 045301 (2022). doi:10.1103/PhysRevE.105.045301
\bibitem{rittig2023} Rittig, J. G., Felton, K. C., Lapkin, A. A. \& Mitsos, A. Gibbs-Duhem-informed neural networks for binary activity coefficient prediction. \textit{Digital Discovery} \textbf{2}, 1752 (2023). doi:10.1039/D3DD00103B
\bibitem{rittig2024} Rittig, J. G. \& Mitsos, A. Thermodynamics-consistent graph neural networks. \textit{Chem. Sci.} \textbf{15}, 18504 (2024). doi:10.1039/D4SC04554H
\bibitem{specht2024} Specht, T. et al. HANNA: hard-constraint neural network for consistent activity coefficient prediction. \textit{Chem. Sci.} \textbf{15}, 19777 (2024). doi:10.1039/D4SC05115G
\bibitem{hoffmann2026} Hoffmann, M. et al. Thermodynamically consistent machine learning model for excess Gibbs energy. \textit{Nat. Commun.} \textbf{17}, 3485 (2026). doi:10.1038/s41467-026-71430-y
\bibitem{cheng2026} Cheng, B. Equivariant learning of a transferable three-dimensional classical density functional. Preprint at https://arxiv.org/abs/2608.13506 (2026).
\bibitem{leenhouts2025} Leenhouts, R. J., Jankelevitch, S., Raike, R., M\"{u}ller, S. \& Vermeire, F. H. Thermodynamics-informed graph neural networks for phase transition enthalpies. \textit{Syst. Control Trans.} \textbf{4}, 1662 (2025). doi:10.69997/sct.140638
\bibitem{kim2016} Kim, H. et al. Characterization of adsorption enthalpy of novel water-stable zeolites and metal-organic frameworks. \textit{Sci. Rep.} \textbf{6}, 19097 (2016). doi:10.1038/srep19097
\bibitem{isodb} Siderius, D. W., Shen, V. K., Johnson III, R. D. \& van Zee, R. D. (eds) \textit{NIST/ARPA-E Database of Novel and Emerging Adsorbent Materials} (National Institute of Standards and Technology, Gaithersburg MD). doi:10.18434/T43882 (records retrieved 13 October 2026)
\bibitem{hanikel2019} Hanikel, N. et al. Rapid cycling and exceptional yield in a metal-organic framework water harvester. \textit{ACS Cent. Sci.} \textbf{5}, 1699 (2019). doi:10.1021/acscentsci.9b00745
\bibitem{taylor2012} Taylor, J. M., Vaidhyanathan, R., Iremonger, S. S. \& Shimizu, G. K. H. Enhancing water stability of metal-organic frameworks via phosphonate monoester linkers. \textit{J. Am. Chem. Soc.} \textbf{134}, 14338 (2012). doi:10.1021/ja306812r
\bibitem{nguyen2022} Nguyen, H. L. et al. Hydrazine-hydrazide-linked covalent organic frameworks for water harvesting. \textit{ACS Cent. Sci.} \textbf{8}, 926 (2022). doi:10.1021/acscentsci.2c00398
\bibitem{schweng2024} Schweng, P., Li, C., Guggenberger, P., Kleitz, F. \& Woodward, R. T. A sulfonated covalent organic framework for atmospheric water harvesting. \textit{ChemSusChem} \textbf{17}, e202301906 (2024). doi:10.1002/cssc.202301906
\bibitem{rieth2017} Rieth, A. J., Yang, S., Wang, E. N. \& Dinc\u{a}, M. Record atmospheric fresh water capture and heat transfer with a material operating at the water uptake reversibility limit. \textit{ACS Cent. Sci.} \textbf{3}, 668 (2017). doi:10.1021/acscentsci.7b00186
\bibitem{tan2021} Tan, K. T., Tao, S., Huang, N. \& Jiang, D. Water cluster in hydrophobic crystalline porous covalent organic frameworks. \textit{Nat. Commun.} \textbf{12}, 6747 (2021). doi:10.1038/s41467-021-27128-4
\bibitem{grunenberg2023} Grunenberg, L. et al. Postsynthetic transformation of imine- into nitrone-linked covalent organic frameworks for atmospheric water harvesting at decreased humidity. \textit{J. Am. Chem. Soc.} \textbf{145}, 13241 (2023). doi:10.1021/jacs.3c02572
\bibitem{lin2006} Lin, X. et al. A porous framework polymer based on a zinc(II) 4,4'-bipyridine-2,6,2',6'-tetracarboxylate: synthesis, structure, and `zeolite-like' behaviors. \textit{J. Am. Chem. Soc.} \textbf{128}, 10745 (2006). doi:10.1021/ja060946u
\bibitem{rezk2012} Rezk, A., Al-Dadah, R., Mahmoud, S. \& Elsayed, A. Characterisation of metal organic frameworks for adsorption cooling. \textit{Int. J. Heat Mass Transf.} \textbf{55}, 7366 (2012). doi:10.1016/j.ijheatmasstransfer.2012.07.068
\bibitem{cadiau2015} Cadiau, A. et al. Design of hydrophilic metal organic framework water adsorbents for heat reallocation. \textit{Adv. Mater.} \textbf{27}, 4775 (2015). doi:10.1002/adma.201502418
\end{thebibliography}


\end{document}
